\section{The cumulative one-loop coefficient and the ultraviolet limit}\label{sec:ol-uv}

This section assembles Sections~\ref{sec:ol-composite}--\ref{sec:ol-feynman} and combines the result with the flow of Part~X. Theorem~\ref{thm:ol-cumulative} is the cumulative one-loop coefficient. Proposition~\ref{prop:flow-F} derives from the flow (Theorem~\ref{thm:sf-flow}) the properties (F1)--(F3) of the couplings that the ultraviolet limit uses, and Theorem~\ref{thm:ol-uv} is our version of the ultraviolet limit, Ba\l aban's Theorem~2 \cite{BalYM4I}, for the couplings $g_k$ of the small-field flow (see the paragraph after it for the comparison with his (0.31)). Ba\l aban's couplings come from his own generations, a different construction (Remark~\ref{rem:sf-flow}(a) and the statement of Proposition~\ref{prop:beta-identification}). The flow reads the coefficients $\beta_j$ off the generations on the given torus (Proposition~\ref{prop:period-iso}), while Part~XI works with marginal tensors in infinite volume. The two are compared through Lemma~\ref{lem:period-vacpol}, with the torus independence of the one-loop parts of the generations (Lemma~\ref{lem:OL}) and the infinite-volume limits of the step functions (Lemma~\ref{lem:ol-limits}); this comparison is (F2). Proposition~\ref{prop:beta-identification} compares $\beta_j$ only with the coefficient that Ba\l aban's definition \cite[(1.20)--(1.22)]{BalYM4I} assigns to the generations of this note, at fixed couplings and as the torus grows; it rests on the generation locality of Section~\ref{sec:gen-locality} and is proved without a rate. Remark~\ref{rem:ol-remains} lists what remains.

\begin{theorem}[The cumulative one-loop coefficient]\label{thm:ol-cumulative}
Let $G$ be compact, connected and simple, with $\beta_\infty=\frac{11C_A}{24\pi^2T_\rho}\log L$ as in Theorem~\ref{thm:ol-log}. Then every $T^{(1)}_{j+1}$ is a multiple $\beta^{(1)}_{j+1}(\delta_{\mu\nu}|p|^2-p_\mu p_\nu)$, and, uniformly in $k$,
\[
\sum_{j<k}T^{(1)}_{j+1}=T_{\Gamma^{(1)}_k}=\big(k\beta_\infty+O(1)\big)\big(\delta_{\mu\nu}|p|^2-p_\mu p_\nu\big),
\]
where $O(1)$ is a number bounded uniformly in $k$.
\end{theorem}

\begin{proof}\hasdetail{thm:ol-cumulative}
The second variations at $1$ of $\Gamma^{(1)}_k$ and of the step functions $\Gamma^{\rm step}_j\circ U_{j+1}$, $j<k$, on the tori converge with uniform exponential decay as the side tends to infinity: for $\Gamma^{(1)}_k$ by Lemma~\ref{lem:ol-torus-limit}(c), and for the step functions by Lemma~\ref{lem:ol-limits} below, whose proof uses the polymer decomposition of Part~X with its parameters; the step functions themselves do not depend on it. The limits are taken along the tori whose side at scale $k$ lies in $\mathfrak N_k$. By Lemma~\ref{lem:torus}(a), with $K=k$, (T0)--(T2) hold at every scale $\le k$, so these tori satisfy the setting of Section~\ref{sec:ol-composite}. For $j<k$ the same torus, read with $K=j+1$, satisfies the Setting of Section~\ref{sec:generations} at the scales $\le j+1$. Its conditions at the scales $i\le j+1$ are among those at the scales $\le k$. The bound on the last side holds because the side at scale $j+1$ is $L^{k-j-1}$ times the side at scale $k$, hence at least $n_{\min}$. So the sides at scale $j+1$ lie in $\mathfrak N_{j+1}$. The limit of the second variation of $\Gamma^{(1)}_k$ (Lemma~\ref{lem:ol-torus-limit}(c)) is the same along this subsequence. So Corollary~\ref{cor:ol-cumulative} applies. By Lemma~\ref{lem:ol-symmetry}(c), these functions have no term linear in $p$, and their marginal tensors are multiples of $\delta_{\mu\nu}|p|^2-p_\mu p_\nu$; the first equality is Corollary~\ref{cor:ol-cumulative}. For the second:
\begin{itemize}
\item by Lemma~\ref{lem:ol-coarse-map}, $T_{\Gamma^{(1)}_k}=T_{\rm fine}$, the $p^2$-tensor of the fine vacuum polarization of $\widehat\Gamma$;
\item by Proposition~\ref{prop:ol-reduction}, $T_{\rm fine}=T_{\rm F}+O(1)$, uniformly in $k$;
\item by Remark~\ref{rem:ol-symmetry}, $T_{\rm F}=\beta_{\rm F}(\delta_{\mu\nu}|p|^2-p_\mu p_\nu)$; by Lemma~\ref{lem:ol-symmetry}, $T_{\Gamma^{(1)}_k}$ is a multiple of the same tensor, so the remainder $O(1)$ is too, and its coefficient is bounded uniformly in $k$;
\item by Theorem~\ref{thm:ol-log}, $\beta_{\rm F}=k\beta_\infty+O(1)$.
\end{itemize}
\end{proof}

\subsection{Averages of linear fields}\label{subsec:ol-linear}

The proof of Lemma~\ref{lem:ol-pullback} uses the averages of the linear fields of Section~\ref{subsec:gen-period}. They are exact.

\begin{lemma}[Averages of linear fields]\label{lem:ol-linear-average}
Let $k\ge1$, and use the units of scale $k$, so that the fine lattice $\eta\Z^4$ has spacing $\eta=L^{-k}$ and the coarse lattice is $\Z^4$. Let $F$ be a curvature datum in $\mathfrak g^c$ and $x_0\in\R^4$, and let $A_F$ and $B_F$ be the linear fields of $F$ with origin $x_0$ (Section~\ref{subsec:gen-period}) on the bonds of $\eta\Z^4$ and of $\Z^4$.
\begin{enumerate}[label=(\alph*)]
\item $\bQ A_F=B_F$, with the gauge-field average $\bQ$ of Definition~\ref{def:averages}, given on $\eta\Z^4$ by the same formula.
\item Let $\bQ^{\rm pl}$ be the linearization at $U=1$ of the average $U\mapsto U^{(k)}$ of Section~\ref{sec:averaging}: $U^{(k)}(c)=e^{\ii t(\bQ^{\rm pl}A)(c)+O(t^2)}$ for $U=e^{\ii\eta tA}$. By Lemma~\ref{lem:one-step-basic}(c), $(\bQ^{\rm pl}A)(c)$ depends only on $A$ on the bonds of $B(c_-)\cup B(c_+)$, so $\bQ^{\rm pl}$ acts on fields on $\eta\Z^4$. Then $\bQ^{\rm pl}A_F=B_F$.
\end{enumerate}
\end{lemma}

This is the input on linear fields of the proof of Lemma~\ref{lem:ol-pullback}: the blocks of Part~IV are centred, and the average of $A_F$ to scale $j+1$ is the linear field of $F$ on the lattice of scale $j+1$, with no gradient ((b) with $j+1$ in place of $k$); in particular the coarse gradient of Lemma~\ref{lem:one-step-expansion}(f) vanishes on linear fields.

\begin{proof}\hasdetail{lem:ol-linear-average}
For a constant field, $A(b)=v_\mu$ on the bonds of direction $\mu$, both $\bQ$ and $\bQ^{\rm pl}$ give the coarse field with the same components: for $\bQ$ since $Q_\mu1=1$ (Lemma~\ref{lem:avg}(b)), for $\bQ^{\rm pl}$ by the gauge covariance below with $\lambda(x)=v\cdot x$. Changing the origin from $x_0$ to $y\in\Z^4$ adds to $A_F$ and to $B_F$ the same constant field, with components $\frac12\sum_\kappa(y-x_0)_\kappa F_{\kappa\mu}$. So it suffices to show, for every coarse point $y$ and direction $\mu$, that the average of $A_F$ with origin $y$ vanishes on $c=(y,y+e_\mu)$; and $B_F(c)=\frac12(m_c-y)^{\sf T}Fe_\mu=\frac14e_\mu^{\sf T}Fe_\mu=0$ for this origin.

(a) For the bonds $b$ of direction $\mu$ starting at $x+i\eta e_\mu$, $0\le i<\eta^{-1}$, $A_F(b)=\frac12\sum_{\kappa\ne\mu}(x-y)_\kappa F_{\kappa\mu}$, since $F_{\mu\mu}=0$. This is linear in $x-y$. The block $B(y)$ is centred at $y$, so $\sum_{x\in B(y)}(x-y)=0$, and Definition~\ref{def:averages} gives $(\bQ A_F)_\mu(y)=0$.

(b) \emph{Covariance.} Let $\lambda$ be a $\mathfrak g$-valued function and $u=e^{-\ii t\lambda}$. Then $(e^{\ii\eta tA})^u=e^{\ii\eta t(A+\partial\lambda)+O(t^2)}$ bondwise, with $\partial\lambda(b)=\eta^{-1}(\lambda(b_+)-\lambda(b_-))$, and Lemma~\ref{lem:one-step-basic}(b) gives $\bQ^{\rm pl}(A+\partial\lambda)=\bQ^{\rm pl}A+\partial'\lambda$, with $\partial'\lambda(c)=\lambda(c_+)-\lambda(c_-)$ on the coarse bonds. For $r\in\mathfrak S_k$, $r(e^{\ii\eta tA})=e^{\ii\eta t\,r_1A}$ with $(r_1A)(b)=A(r^{-1}b)$ and $A(-b)=-A(b)$, and Lemma~\ref{lem:one-step-basic}(d) gives $\bQ^{\rm pl}r_1=r_1\bQ^{\rm pl}$. These identities are stated on the torus; by the locality, each of them at a coarse bond involves the fields only on the blocks of that bond and of its image, so they hold on $\eta\Z^4$ for all fields, after the fields are replaced near these blocks by periodic fields on a large torus (for the second identity, by an $r$-invariant cut-off). Both extend to complex fields by linearity.

\emph{Reflections.} Let $\nu\ne\mu$, and let $r$ be the reflection of the $\nu$-th coordinate about $y$, an element of $\mathfrak S_k$ with linear part $R_\nu$. It maps $c$ onto itself, with its orientation, and $r_1A_F=A_{R_\nu\cdot F}$, with origin $ry=y$, by \eqref{eq:linear-field-covariance}. So $\ell(F)=(\bQ^{\rm pl}A_F)(c)$ satisfies $\ell(R_\nu\cdot F)=(r_1\bQ^{\rm pl}A_F)(c)=(\bQ^{\rm pl}A_F)(r^{-1}c)=\ell(F)$. Write $\ell(F)=\sum_{\kappa<\lambda}\ell_{\kappa\lambda}(F_{\kappa\lambda})$, with linear maps $\ell_{\kappa\lambda}$ on $\mathfrak g^c$. The map $F\mapsto R_\nu\cdot F$ changes the sign of $F_{\kappa\lambda}$ if exactly one of $\kappa,\lambda$ equals $\nu$, and keeps the other entries. Every pair $\kappa<\lambda$ contains an index $\nu\ne\mu$, and the other index of the pair is not $\nu$. So $\ell_{\kappa\lambda}=-\ell_{\kappa\lambda}$ for every pair, and $\ell=0$.
\end{proof}

\subsection{The coefficients of the flow}\label{subsec:ol-flow-coeff}

Let the parameters and the torus be as in Section~\ref{sec:flow4} (Setting), with $n_K\ge n_{\min}$, and let $\gamma$ satisfy ($\Gamma$1)--($\Gamma$3) of Remark~\ref{rem:flow-constants}. For $\pi=(g_0,\dots,g_j)\in\Sigma_j$ (Section~\ref{subsec:flow4-flow}) the generations $E^{(i)}=E^{(i)}[\pi]$ and the coefficients $\beta_i=\beta_i[\pi]$, $i\le j+1$, are those of Proposition~\ref{prop:flow-continuity}; so $\beta_{j+1}$ is a function of $(g_0,\dots,g_j)$. The couplings of the flow are those of Corollary~\ref{cor:bare-continuity}: $g_{j+1}^{-2}=g_j^{-2}-\beta_{j+1}$ as long as $g_0,\dots,g_j\in(0,\gamma]$, and then $g_{j+1}\le2g_j$. In this section the coupling $g$ is that of the last scale, $g_K=g$, and the bare coupling is $g_0$; in Sections~\ref{sec:rg-step} and~\ref{sec:target-flow}, by contrast, $g=g_0$ denotes the bare coupling, and in Sections~\ref{sec:ol-composite}--\ref{sec:ol-feynman} $g$ is the single coupling of the Laplace asymptotics.

\paragraph{What Part~X provides.} The density of every scale has the generation form of Remark~\ref{rem:flow-status}. A generation of scale $j$ is a family of normalized, gauge-invariant graded activities on the polymers of scale $j$, defined on the domains $\mathcal U^\flat(X)$ with $\delta_1^\flat=\frac14\delta_1$, with the properties (G1)--(G4) and the norm $\|E\|_j=\sup_Xe^{\kappa_nd_j(X)}\|E(X)\|$ (Definition~\ref{def:generation}). A step produces a raw generation $E'$ (Definition~\ref{def:raw-new}); its average over the symmetries of the grid is the new generation (Definition~\ref{def:symmetrization}). The coefficient of a generation is read off from its period density, the quadratic part of its activities at the linear fields $A_F$ of constant curvature, per grid cube: $\bar T_E=\beta_E\bar T_S$ (Definition~\ref{def:period-density}, Proposition~\ref{prop:period-iso}). Along the flow, $\|E^{(j)}\|_j\le\mathsf K$ and $|\beta_j|\le C_\beta\mathsf K$ (Theorem~\ref{thm:sf-flow}). As in Section~\ref{sec:flow}, $\gamma$ is the bound on the couplings, not a decay rate; ($\Gamma$1)--($\Gamma$4) are the conditions on it of Remark~\ref{rem:flow-constants}.

\paragraph{The volume coefficient.} Let $E$ be a family of normalized graded activities at the polymers of scale $j$, on domains containing the $\mathcal U^\flat(X)$; for instance a family as in Definition~\ref{def:period-density}, or a raw generation. We do not require it to be real. Put $\|E\|=\sup_Xe^{\kappa_nd_j(X)}\sup_{\mathcal U^\flat(X)}|E(X,\cdot)|$, and, with $(T_a)$ and $F^{\mu\nu,a}$ as in Proposition~\ref{prop:period-iso}(c),
\begin{equation}\label{eq:beta-vol}
\bar T^{\rm vol}_E(F)=n_j^{-4}\sum_Xc^X(F),\qquad\beta^{\rm vol}_E=\frac1{3\dim\mathfrak g}\sum_{\mu<\nu}\ \sum_a\bar T^{\rm vol}_E\big(F^{\mu\nu,a}\big),
\end{equation}
the first sum running over the polymers that are not wrapping. The definition of $c^X$ and Lemma~\ref{lem:period-density}(a), (c) use neither the reality of $E$ nor (G2): the bound of (c) holds for each grid cube $Q$, and $\bar T^{\rm vol}_E(F)=n_j^{-4}\sum_Q\sum_{X\supset Q}c^X(F)/|X|$ is an average over the $Q$. So $\beta^{\rm vol}_E$ is linear in $E$, and, as in the proof of Proposition~\ref{prop:period-iso}(b),
\begin{equation}\label{eq:beta-vol-bound}
|\beta^{\rm vol}_E|\le C_\beta\|E\|,
\end{equation}
with $C_\beta$ of Proposition~\ref{prop:period-iso}(b). If $E$ has (G2) and (G4), then $\beta^{\rm vol}_E=\beta_E$, by Lemma~\ref{lem:period-density}(c) and Proposition~\ref{prop:period-iso}(a). If $E'$ is a raw generation, then $\beta_{E'^{\rm sym}}=\beta^{\rm vol}_{E'}$, by Proposition~\ref{prop:period-iso}(c).

\paragraph{The one-loop part.} Let $E'$ be the raw new generation of the step from scale $k$ to scale $k+1$ (Definition~\ref{def:raw-new}). Write $E'=E'_{\rm ol}+E'_\#$, with
\[
\begin{aligned}
E'_{\rm ol}(Y',U)&=\sum_{X:\,\bar X=Y'}\mathfrak e(X,U)+\sum_{c:\,Y'_c=Y'}\log\mathfrak j_c(U),\\
E'_\#(Y',U)&=\sum_{Y:\,\hat Y=Y'}\big(E^\#(Y,U)-E^\#(Y,1)\big),
\end{aligned}
\]
the sums as in Definition~\ref{def:raw-new}.
The one-loop part $E'_{\rm ol}$ involves neither the coupling, nor the old activities, nor the white-noise cutoff: it is a function of the reference alone, through the fluctuation operator and the base (Propositions~\ref{prop:one-loop-polymer} and~\ref{prop:jacobians}; Proposition~\ref{prop:first-step}(c) at $k=0$). The proof of Lemma~\ref{lem:raw-new}(a), (b), restricted to its first two terms, shows that $E'_{\rm ol}$ has (G1) on the domains $\mathcal U(Y')$ of scale $k+1$, with $\|E'_{\rm ol}\|^{\rm raw}_{k+1}\le\mathsf K$. The \emph{one-loop part} of $\beta_{k+1}$ is
\[
\beta^{\rm ol}_{k+1}=\beta^{\rm vol}_{E'_{\rm ol}}.
\]
It depends on $k$, on the torus and on the parameters fixed before $\gamma$, not on the couplings. Since $\beta_{k+1}=\beta^{\rm vol}_{E'}$ and $\beta^{\rm vol}$ is linear, $\beta_{k+1}=\beta^{\rm ol}_{k+1}+\beta^{\rm vol}_{E'_\#}$.

\paragraph{The properties of the flow used in the ultraviolet limit.} With $\beta^{(1)}_{j+1}$ the coefficient of $T^{(1)}_{j+1}$ (Theorem~\ref{thm:ol-cumulative}):
\begin{enumerate}[label=(F\arabic*)]
\item If $g_0,\dots,g_k\in(0,\gamma]$, the densities $\rho_1,\dots,\rho_k$ of the flow have the generation form (Definition~\ref{def:gen-form}), with $g_{j+1}^{-2}=g_j^{-2}-\beta_{j+1}$. The coefficient $\beta_{j+1}$ is a continuous function of $(g_0,\dots,g_j)$ on $\Sigma_j$, and $|\beta_{j+1}|\le C_\beta\mathsf K$.
\item $|\beta^{\rm ol}_{j+1}-\beta^{(1)}_{j+1}|\le C_2e^{-c_2n_{j+1}/M}$ for $j<K$, with $C_2,c_2>0$ independent of $j$, $K$ and the torus.
\item $|\beta_{j+1}-\beta^{\rm ol}_{j+1}|\le C_3g_j^2p_0(g_j)^6$, uniformly in $j$ and in the earlier couplings, for $(g_0,\dots,g_j)\in\Sigma_j$. Here $p_0$ is the white-noise cutoff of Section~\ref{sec:flow}.
\end{enumerate}
In Ba\l aban's formulation of (F2), the marginal parts of the old generations are cancelled before the one-loop part is read off. Here the cancellation is built into the step (Remark~\ref{rem:flow-F}(a)).

\subsection{Period density and vacuum polarization}\label{subsec:ol-period-vacpol}

\paragraph{Flat domains and the infinite lattice.} Use the units of scale $j$. For a polymer $X$ of scale $j$ that is not wrapping and a field $B$ on the bonds of a lift of $X^{(3)}$, let $V_B$ be the configuration equal to $e^{\ii\eta B}$ on the bonds of $X^{(3)}$, read in the lift, and to $1$ elsewhere, as in the proof of Lemma~\ref{lem:period-density}(a). The \emph{flat domain} $\mathcal U^\flat_1(X)$ is the set of the $V_B$ with $|B|<\delta_0$ and $|\nabla B|+[\nabla B]_\alpha<\delta_1^\flat$ on $X^{(3)}$, the differences and quotients taken with the real part $1$. By Definition~\ref{def:graded} with the real part $1$, $\mathcal U^\flat_1(X)\subset\mathcal U^\flat(X)$. The flat domain is defined by conditions on the bonds of $X^{(3)}$ alone, so it is the same set on every torus on which $X$ does not wrap. The numbers $c^X(F)$ and the second variations at $1$ use $E(X,\cdot)$ only on $\mathcal U^\flat_1(X)$ (proof of Lemma~\ref{lem:period-density}(a)).

A \emph{polymer of scale $j$ of the infinite lattice} is a finite $*$-connected union of cubes of side $M$ centred at the points of $M\Z^4$, read on $\eta\Z^4$, $\eta=L^{-j}$. The tree length $d_j$, the number $|X|$ of cubes, $X^{(3)}$, the flat domains and \eqref{eq:dM-facts} are as on the torus, and these polymers count as not wrapping. A family $E^\infty$ of \emph{flat activities} indexed by them consists of holomorphic functions $E^\infty(X,\cdot)$ on $\mathcal U^\flat_1(X)$ that depend only on $B$ on the bonds of $X$, with $E^\infty(X,1)=0$, real for real $B$, and invariant under small real gauge transformations: $E^\infty(X,V_{B'})=E^\infty(X,V_B)$ whenever $V_B,V_{B'}\in\mathcal U^\flat_1(X)$ and $e^{\ii\eta B'(b)}=e^{\ii\lambda(b_-)}e^{\ii\eta B(b)}e^{-\ii\lambda(b_+)}$ with a small real $\lambda$. For such a family with (G2) and $\|E^\infty\|=\sup_Xe^{\kappa_nd_j(X)}\sup_{\mathcal U^\flat_1(X)}|E^\infty(X,\cdot)|<\infty$, the numbers $c^X(F)$ are defined as in Definition~\ref{def:period-density}, and we put, for a grid cube $Q$,
\[
\bar T_{E^\infty}(F)=M^{-4}\sum_{X\supset Q}\frac{c^X(F)}{|X|},\qquad\beta^\infty_{E^\infty}=\frac1{3\dim\mathfrak g}\sum_{\mu<\nu}\ \sum_a\bar T_{E^\infty}\big(F^{\mu\nu,a}\big).
\]
By (G2) and the proofs of Lemma~\ref{lem:period-density}(a), (c), which use only these properties, the sums converge absolutely and do not depend on $Q$. Let $\widetilde\Pi$ be the kernel, in the trace pairing $\sum_{b,b'}\widetilde\Pi(b,b')\operatorname{tr}(A(b)A(b'))$, of the second variation at $A=0$ of $A\mapsto\sum_XE^\infty(X,e^{\ii\eta A})$, for small fields $A$ of finite support; only the polymers that contain a bond of the support of $A$ contribute, since $E^\infty(X,\cdot)$ depends only on the bonds of $X$ and $E^\infty(X,1)=0$, and the sum converges absolutely by \eqref{eq:dM-facts}. The kernel is invariant under the translations by $M\Z^4$, and $\widetilde\Pi^{\rm av}=M^{-4}\sum_a\tau_a\widetilde\Pi\tau_{-a}$, the average over $a\in\{0,\dots,M-1\}^4$, is invariant under those by $\Z^4$.

\paragraph{Torus-local families.} Let $\mathfrak N_j$ be the set of the sides $\tilde n$ for which the torus whose lattice of scale $j$ has side $\tilde n$ satisfies the hypotheses of the Setting of Section~\ref{sec:generations} at the scales $\le j$ (with $K=j$). A collection $(E[\tilde n])_{\tilde n\in\mathfrak N_j}$, where $E[\tilde n]$ is a family of normalized graded activities at the polymers of scale $j$ of the torus of side $\tilde n$, on domains containing the $\mathcal U^\flat(X)$, real at real configurations, with $\sup_{\tilde n}\|E[\tilde n]\|<\infty$, is \emph{torus-local} if there are functions $E^\infty(X,\cdot)$ on the flat domains of the polymers $X$ of scale $j$ of the infinite lattice, with $E^\infty(X+a,\tau_aU)=E^\infty(X,U)$ for $a\in M\Z^4$, and constants $C,c>0$, such that
\begin{equation}\label{eq:torus-local}
\big|E[\tilde n](X,U)-E^\infty(\tilde X,\tilde U)\big|\le Ce^{-c\tilde n/M}e^{-\kappa_nd_j(X)}
\end{equation}
for every $\tilde n\in\mathfrak N_j$, every polymer $X$ of the torus of side $\tilde n$ that is not wrapping, every lift $\tilde X$ of $X$ (Section~\ref{subsec:gen-period}) and every $U\in\mathcal U^\flat_1(X)$, with $\tilde U\in\mathcal U^\flat_1(\tilde X)$ its lift. We call $E^\infty$ its \emph{limit}.

\paragraph{One-loop locality.} The one-loop parts of the raw new generations form torus-local families, with no error term. This rests on the torus independence of the one-loop pieces and of the Jacobian terms (Section~\ref{subsec:torus-independence}).

\begin{lemma}[One-loop locality]\label{lem:OL}
Let $k\ge0$. For $\tilde n\in\mathfrak N_{k+1}$ let $E'_{\rm ol}[\tilde n]$ be the one-loop part of the raw new generation of the step from scale $k$ to scale $k+1$ on the torus whose lattice of scale $k+1$ has side $\tilde n$. Then the collection $(E'_{\rm ol}[\tilde n])_{\tilde n\in\mathfrak N_{k+1}}$ is torus-local, and $\|E'_{\rm ol}[\tilde n]\|\le\mathsf K$ for every $\tilde n$. Its limit $E^\infty$ satisfies
\[
E'_{\rm ol}[\tilde n](Y',U)=E^\infty(\tilde Y',\tilde U)
\]
for every $\tilde n\in\mathfrak N_{k+1}$, every polymer $Y'$ of scale $k+1$ of the torus of side $\tilde n$ that is not wrapping, every lift $\tilde Y'$ of $Y'$ and every $U\in\mathcal U^\flat_1(Y')$, with its lift $\tilde U$. So the left side of \eqref{eq:torus-local} vanishes, and the bound holds for all $C,c>0$.
\end{lemma}

\begin{proof}\hasdetail{lem:OL}
Let $\tilde n\in\mathfrak N_{k+1}$. By Lemma~\ref{lem:torus}(a), the torus satisfies (T0)--(T2) at step $k$, so its side at scale $k$, which is $L\tilde n$, is admissible in the sense of Section~\ref{subsec:torus-independence}, and Proposition~\ref{prop:one-loop-local} applies to it (with the pieces of Proposition~\ref{prop:first-step}(c) if $k=0$). The activities $E'_{\rm ol}[\tilde n](Y',\cdot)$ are normalized and real at real configurations, on the domains $\mathcal U(Y')\supset\mathcal U^\flat(Y')$ of scale $k+1$, and $\|E'_{\rm ol}[\tilde n]\|\le\|E'_{\rm ol}[\tilde n]\|^{\rm raw}_{k+1}\le\mathsf K$ (paragraph ``The one-loop part'' of Section~\ref{subsec:ol-flow-coeff}), with $\mathsf K$ independent of the torus.

The closure $\bar X=\operatorname{cl}_{k+1}X$ (Section~\ref{subsec:flow-reblock}) and the polymer $Y'_c$, the union of the grid cubes of scale $k+1$ that meet $B'(c_-)\cup B'(c_+)$ (the paragraph before Definition~\ref{def:raw-new}), are given by rules that are local in the grids of the scales $k$ and $k+1$. So they are defined on the infinite lattice, and on sets that do not wrap they commute with the projection. For a polymer $\tilde Y'$ of scale $k+1$ of the infinite lattice and $\tilde U=V_B\in\mathcal U^\flat_1(\tilde Y')$ let $E^\infty(\tilde Y',\tilde U)$ be the sum of the values $E^\infty_k(\tilde X,\cdot)$ over the polymers $\tilde X$ of scale $k$ of the infinite lattice with $\operatorname{cl}_{k+1}\tilde X=\tilde Y'$, and of the values $J^\infty_k(\tilde c,\cdot)$ over the positively oriented bonds $\tilde c$ of $L\Z^4$ with $\tilde Y'_{\tilde c}=\tilde Y'$, at the restrictions of the field $B$, read in the units of step $k$ (Proposition~\ref{prop:one-loop-local}(d)). These sums are finite, since $\tilde X$ and $P_{\tilde c}$ lie in $\tilde Y'$, and the restrictions lie in the domains of $E^\infty_k$ and $J^\infty_k$, by Remark~\ref{rem:one-loop-local-use}(a). Let $Y'$ be a polymer of scale $k+1$ of the torus that is not wrapping, and $\tilde Y'$ a lift. The polymers $X$ with $\bar X=Y'$ and the sets $P_c$ of Lemma~\ref{lem:local-inventory}(d) with $Y'_c=Y'$ lie in $Y'$, so lifting into $\tilde Y'$ is a bijection from them onto the $\tilde X$ and $\tilde c$ of the sums above. By Proposition~\ref{prop:one-loop-local}(d), each term of $E'_{\rm ol}[\tilde n](Y',U)$ is the value of $E^\infty_k$ or $J^\infty_k$ at the lift, so $E'_{\rm ol}[\tilde n](Y',U)=E^\infty(\tilde Y',\tilde U)$ (Remark~\ref{rem:one-loop-local-use}(a)). Finally, a translation by $a\in M\Z^4$ in the units of scale $k+1$ is a translation by $La\in M\Z^4$ in those of step $k$; it commutes with the closure and with $c\mapsto Y'_c$, so the grid covariance of $E^\infty_k$ and $J^\infty_k$ in Proposition~\ref{prop:one-loop-local}(d) gives $E^\infty(\tilde Y'+a,\tau_a\tilde U)=E^\infty(\tilde Y',\tilde U)$.
\end{proof}

For a bond $b$ let $m_b$ be its midpoint and $\mu(b)$ its direction.

\begin{lemma}[Period density and vacuum polarization]\label{lem:period-vacpol}
Let $G$ be compact, connected and simple, let $1\le j\le K$, and use the units of scale $j$.
\begin{enumerate}[label=(\alph*)]
\item \emph{(A single polymer.)} Let $E$ be a family as in Definition~\ref{def:period-density}, and $X$ a polymer that is not wrapping. There is a unique real symmetric kernel $\kappa_X$ on the positively oriented bonds of a lift of $X$ such that $\frac12\frac{d^2}{dt^2}E(X,V_{tB})|_{t=0}=\sum_{b,b'}\kappa_X(b,b')\operatorname{tr}(B(b)B(b'))$ for all $\mathfrak g^c$-valued $B$. It satisfies $\sum_{b'}\kappa_X(b,b')\,\partial\varphi(b')=0$ for every scalar function $\varphi$ and every $b$. Put, for $p\in\R^4$,
\[
m_X(p)_{\mu\nu}=-\frac14\sum_{b,b'}\kappa_X(b,b')\big(p\cdot(m_b-m_{b'})\big)^2\big(\delta_{\mu(b)\mu}\delta_{\mu(b')\nu}+\delta_{\mu(b)\nu}\delta_{\mu(b')\mu}\big).
\]
Then, for $p,a\in\R^4$ and $T\in\mathfrak g^c$, with the curvature datum $(p\wedge a\otimes T)_{\mu\nu}=(p_\mu a_\nu-p_\nu a_\mu)T$,
\[
c^X(p\wedge a\otimes T)=a^{\sf T}m_X(p)\,a\,\operatorname{tr}(T^2).
\]
\item \emph{(The fine tensor.)} Put $T^{\rm f}_E(p)=2n_j^{-4}\sum_Xm_X(p)$, the sum over the polymers that are not wrapping. Then:
\begin{enumerate}[label=(\roman*)]
\item $\bar T^{\rm vol}_E(p\wedge a\otimes T)=\frac12a^{\sf T}T^{\rm f}_E(p)\,a\,\operatorname{tr}(T^2)$;
\item the isotropic coefficient of $T^{\rm f}_E$, in the sense of Definition~\ref{def:marginal-tensor}, is $\beta^{\rm vol}_E$.
\end{enumerate}
\item \emph{(The coarse tensor.)} Let $E^\infty$ be a family of flat activities of scale $j$ of the infinite lattice with (G2) and $\|E^\infty\|<\infty$, and put $T^{\rm f}_{E^\infty}(p)=2M^{-4}\sum_{X\supset Q}m_X(p)/|X|$. Then (a) and (b) hold with $\bar T_{E^\infty}$, $T^{\rm f}_{E^\infty}$, $\beta^\infty_{E^\infty}$ in place of $\bar T^{\rm vol}_E$, $T^{\rm f}_E$, $\beta^{\rm vol}_E$, and $T^{\rm f}_{E^\infty}(p)$ is the coefficient of order $|p|^2$ of $\widetilde\Pi^{\rm av}(p,p)$, the diagonal Bloch component per unit volume in the bond-midpoint convention. Let $\bH_\infty$ be the infinite-volume minimizer map of Proposition~\ref{prop:ol-IM}, and suppose that $\Pi_\Phi=\bH_\infty^{\sf T}\widetilde\Pi\bH_\infty$ is invariant under the translations of $\Z^4$ and under the point reflection $x\mapsto-x$. Then $\Pi_\Phi(p)=O(|p|^2)$, and the marginal tensor of $\Pi_\Phi$ (Definition~\ref{def:marginal-tensor}) is $T^{\rm f}_{E^\infty}$. In particular its isotropic coefficient is $\beta^\infty_{E^\infty}$.
\item \emph{(Torus-local families.)} Let $(E[\tilde n])_{\tilde n\in\mathfrak N_j}$ be torus-local, with limit $E^\infty$, and $\mathsf K'=\sup_{\tilde n}\|E[\tilde n]\|$. Then $E^\infty$ is a family of flat activities with (G2) and $\|E^\infty\|\le\mathsf K'$, and there are $C',c'>0$, depending only on $d$, $M$, $\delta_1$, $\mathsf K'$ and the constants of \eqref{eq:torus-local}, with
\[
\big|\beta^{\rm vol}_{E[\tilde n]}-\beta^\infty_{E^\infty}\big|\le C'e^{-c'\tilde n/M}\qquad(\tilde n\in\mathfrak N_j).
\]
Moreover, the kernels of the second variations at $A=0$ of $A\mapsto\sum_XE[\tilde n](X,e^{\ii\eta A})$ on the tori converge to $\widetilde\Pi$ with uniform decay, in the sense of Lemma~\ref{lem:ol-folding}(e): entrywise as $\tilde n\to\infty$, and with a bound $C_\eta e^{-\gamma|m_b-m_{b'}|_{\tilde n}/M}$, uniform in $\tilde n$, where $|\cdot|_{\tilde n}$ is the distance on the torus and $C_\eta$, $\gamma>0$ may depend on $\eta$.
\end{enumerate}
\end{lemma}

The tensor $T^{\rm f}_E$ is the $p^2$-tensor of the total second variation of the pieces of $E$ in the fine field, before the composition with the background field. By (b) it determines the isotropic part $\beta^{\rm vol}_E$ of the period density, and $\bar T^{\rm vol}_E$ on the data $p\wedge a\otimes T$. Since a quadratic form on $\Lambda^2\R^4$ that vanishes on the decomposable bivectors is a multiple of the form of $\operatorname{Pf}$ (the Pl\"ucker relation), it determines $\bar T^{\rm vol}_E$ up to a multiple of $\operatorname{Pf}$, the part that the symmetrization removes (Remark~\ref{rem:period-density}(a)). So the caveat of Remark~\ref{rem:period-density}(c), that $\beta_E$ is a property of the family, does not matter for the one-loop part: $\beta^{\rm ol}_{k+1}$ is the isotropic part of $\bar T^{\rm vol}_{E'_{\rm ol}}$ (Proposition~\ref{prop:period-iso}(c)), determined by the fine tensor of the one-loop pieces, and by Lemma~\ref{lem:OL} and (d) it is, up to $O(e^{-n_{k+1}/(4M)})$, the coefficient $\beta^\infty_{E^\infty}$ of the infinite-volume limit $E^\infty$ of the one-loop parts. For the plaquette family of Lemma~\ref{lem:period-density}(d), $\bar T^{\rm vol}=\bar T_S=\frac12\operatorname{Iso}$, and (b)(i) gives $T^{\rm f}(p)_{\mu\nu}=\delta_{\mu\nu}|p|^2-p_\mu p_\nu$, the marginal tensor of $S$ in the trace pairing (proof of Lemma~\ref{lem:ol-classical}). So the coefficients $\beta_E$, defined relative to $\bar T_S$, and the isotropic coefficients of Part~XI are normalized alike, with conversion factor $1$.

\begin{proof}\hasdetail{lem:period-vacpol}
(a) As in the proof of Lemma~\ref{lem:period-density}(a), $\Phi_X(B)=E(X,V_B)$ is holomorphic for small $B$, depends only on $B$ on the bonds of $X$, and $d\Phi_X(0)=0$. For constant $v\in G$ near $1$, $V_{vBv^{-1}}=(V_B)^v$ on the bonds of $X^{(3)}$, so $\frac12d^2\Phi_X(0)$ is a quadratic form in $(B(b))_b$ invariant under the simultaneous adjoint action of these $v$, hence of $G$, which is connected. Each block of its polar form, $(T,T')\mapsto k_{bb'}(T,T')$, is then an invariant bilinear form on $\mathfrak g^c$, and, as in the proof of Lemma~\ref{lem:b4-forms}, $k_{bb'}(T,T')=\kappa_X(b,b')\operatorname{tr}(TT')$, with $\kappa_X$ symmetric, and real since $E$ is real at real configurations and $\operatorname{tr}(T^2)>0$ for $0\ne T\in\mathfrak g$. The proof of Lemma~\ref{lem:period-density}(a) gives $d^2\Phi_X(0)[B,\partial\lambda]=0$ for all $B$ and $\lambda$; with $B$ supported on one bond and $\lambda=\varphi T$, this is the second identity.

Let $F=p\wedge a\otimes T$, and take the origin $0$ in the lift. On a positively oriented bond $b$ of direction $\nu$, $A_F(b)=\frac12\sum_\kappa(m_b)_\kappa(p_\kappa a_\nu-p_\nu a_\kappa)T=\frac12\big((p\cdot m_b)a_\nu-(a\cdot m_b)p_\nu\big)T$. Let $\varphi(x)=\frac12(p\cdot x)(a\cdot x)$. Since $b_\pm=m_b\pm\frac12\eta e_\nu$,
\[
\partial\varphi(b)=\eta^{-1}\big(\varphi(b_+)-\varphi(b_-)\big)=\tfrac12\big((p\cdot m_b)a_\nu+(a\cdot m_b)p_\nu\big)
\]
exactly. Hence $A_F=(B_1-\partial\varphi)T$, with $B_1(b)=(p\cdot m_b)a_{\mu(b)}$. Write $\kappa_X(B,B')=\sum_{b,b'}\kappa_X(b,b')B(b)B'(b')$ for scalar fields. By the second identity, $c^X(F)=\frac12d^2\Phi_X(0)[A_F,A_F]=\kappa_X(B_1,B_1)\operatorname{tr}(T^2)$. Let $B_0(b)=a_{\mu(b)}=\partial\psi(b)$ with $\psi(x)=a\cdot x$, and $B_2(b)=(p\cdot m_b)^2a_{\mu(b)}$. Expanding the square,
\[
\begin{aligned}
\sum_{b,b'}\kappa_X(b,b')\,a_{\mu(b)}a_{\mu(b')}\big(p\cdot m_b-p\cdot m_{b'}\big)^2&=\kappa_X(B_2,B_0)-2\kappa_X(B_1,B_1)+\kappa_X(B_0,B_2)\\
&=-2\kappa_X(B_1,B_1),
\end{aligned}
\]
by the second identity and the symmetry of $\kappa_X$. By the definition of $m_X$, the left side is $-2a^{\sf T}m_X(p)a$. So $a^{\sf T}m_X(p)a=\kappa_X(B_1,B_1)$, which gives (a). The kernel and $m_X$ do not depend on the lift, since only differences of midpoints enter.

(b) (i) Sum (a) over $X$. For (ii): by Lemma~\ref{lem:period-density}(a), $\bar T^{\rm vol}_E$ is a quadratic form, invariant under conjugation, so, by the first part of the proof of Lemma~\ref{lem:b4-forms}, $\bar T^{\rm vol}_E(F)=\sum_{\mu<\nu}\sum_{\kappa<\lambda}Q(e_{\mu\nu},e_{\kappa\lambda})\operatorname{tr}(F_{\mu\nu}F_{\kappa\lambda})$ with a symmetric bilinear form $Q$ on $\Lambda^2\C^4$; for $w\in\Lambda^2\R^4$, $\bar T^{\rm vol}_E(w\otimes T)=Q(w,w)\operatorname{tr}(T^2)$. By (i), $Q(p\wedge a,p\wedge a)=\frac12a^{\sf T}T^{\rm f}_E(p)a$ for real $p,a$, and, polarizing in $a$,
\begin{equation}\label{eq:Tf-Q}
T^{\rm f}_E(p)_{\mu\nu}=2Q(p\wedge e_\mu,p\wedge e_\nu).
\end{equation}

(ii) For $R\in B_4$ put $Q^R(w,w')=Q(Rw,Rw')$, with $R$ acting on $\Lambda^2$, and $T^R(p)=R^{\sf T}T(Rp)R$ for $p^2$-tensors $T$ (in this part of the proof $T$ is a tensor, not the element $T\in\mathfrak g^c$ of the statement). By \eqref{eq:Tf-Q}, the tensor built from $Q^R$ is $(T^{\rm f}_E)^R$, since $2Q(Rp\wedge Re_\mu,Rp\wedge Re_\nu)=\sum_{\alpha,\beta}R_{\alpha\mu}R_{\beta\nu}\,2Q(Rp\wedge e_\alpha,Rp\wedge e_\beta)$. The map $T\mapsto T^R$ is the action of $R$ on the indices of the coefficient tensor, which is orthogonal for the Euclidean inner product, and it fixes the tensor of $\delta_{\mu\nu}|p|^2-p_\mu p_\nu$. So the isotropic coefficient of $T^{\rm f}_E$ is that of the tensor built from $\bar Q=|B_4|^{-1}\sum_RQ^R$. The form $\bar Q$ defines a quadratic form on the curvature data that is invariant under $B_4$ and under conjugation, so $\bar Q=\beta'I$ by Lemma~\ref{lem:b4-forms}(a), with $I(w,w')=\sum_{\mu<\nu}w_{\mu\nu}w'_{\mu\nu}$, the form of Iso. Its tensor is $2\beta'I(p\wedge e_\mu,p\wedge e_\nu)=2\beta'(\delta_{\mu\nu}|p|^2-p_\mu p_\nu)$, so the isotropic coefficient of $T^{\rm f}_E$ is $2\beta'$. On the other hand, every $R\in B_4$ permutes the $e_{\mu\nu}$ up to signs, so $\sum_{\mu<\nu}Q(e_{\mu\nu},e_{\mu\nu})=\sum_{\mu<\nu}\bar Q(e_{\mu\nu},e_{\mu\nu})=6\beta'$, and, with $\operatorname{tr}(T_a^2)=1$, $\beta^{\rm vol}_E=(3\dim\mathfrak g)^{-1}\dim\mathfrak g\cdot6\beta'=2\beta'$.

(c) Part (a) concerns one polymer and holds verbatim, since its proof uses $E(X,\cdot)$ only on the flat domain and its invariance under small gauge transformations. Part (b) holds with the cell sums: (i) follows from (a), summed with the weights $M^{-4}|X|^{-1}$, and (ii) is a statement about the forms $\bar T_{E^\infty}$ and $T^{\rm f}_{E^\infty}$ alone. The sums converge absolutely: by Cauchy's estimate, as in Lemma~\ref{lem:period-density}(a), $|\kappa_X(b,b')|\le C_\eta\|E^\infty\|e^{-\kappa_nd_j(X)}$ with $C_\eta$ depending on $\eta$ and the domain constants, a polymer that contains $b$ and $b'$ has $d_j(X)\ge|m_b-m_{b'}|/M-2$ (as in Lemma~\ref{lem:torus}(b)), and \eqref{eq:dM-facts} applies. In particular $\widetilde\Pi=2\sum_X\kappa_X$ has exponential decay. The diagonal Bloch component per unit volume is $\widetilde\Pi^{\rm av}(p,p)_{\mu\nu}=M^{-4}\sum_{b,b'}e^{-\ii p\cdot(m_b-m_{b'})}\widetilde\Pi(b,b')$, the sum over the bonds $b$ of direction $\mu$ with midpoints in an $M$-cell and the bonds $b'$ of direction $\nu$. By (G2), every orbit of the translations by $M\Z^4$ on the polymers contains exactly $|X|$ polymers that contain $Q$. Hence
\[
\widetilde\Pi^{\rm av}(p,p)_{\mu\nu}=2M^{-4}\sum_{X\supset Q}\frac1{|X|}\sum_{b,b'}e^{-\ii p\cdot(m_b-m_{b'})}\kappa_X(b,b')\,\delta_{\mu(b)\mu}\delta_{\mu(b')\nu},
\]
the inner sum over all pairs of bonds of $X$.
Expand the exponential. The term of order $0$ vanishes by the second identity of (a), since a constant field is a gradient. The term of order $1$ is antisymmetric in $\mu,\nu$, by the symmetry of $\kappa_X$. The term of order $2$ is symmetric in $\mu,\nu$ and equals $T^{\rm f}_{E^\infty}(p)$. So $\widetilde\Pi^{\rm av}(p,p)=\mathsf l(p)+T^{\rm f}_{E^\infty}(p)+O(|p|^3)$, with $\mathsf l$ linear and antisymmetric.

The operator $\bH_\infty$ lies in $\mathscr P$ (Proposition~\ref{prop:ol-IM}(a)), so it commutes with the translations of $\Z^4$; since $\Pi_\Phi$ is invariant under them, $\bH_\infty^{\sf T}\widetilde\Pi^{\rm av}\bH_\infty=\Pi_\Phi$. The kernel $\widetilde\Pi^{\rm av}$ is real and symmetric, since the $\kappa_X$ are, and it is invariant under the translations of $\Z^4$, with exponential decay. It annihilates the constant fields, by the second identity of (a), since a constant field is the gradient of a linear function. So Lemma~\ref{lem:ol-coarse-map}, with $j$ in place of $k$, applies to $\widetilde\Pi^{\rm av}$. Its first part gives $\Pi_\Phi(p)=\widetilde\Pi^{\rm av}(p,p)+O(|p|^3)$; its proof uses Proposition~\ref{prop:ol-IM}(e) and Corollary~\ref{cor:ol-umklapp}, and the Umklapp sum is finite. The point reflection gives $\Pi_\Phi(-p)=\Pi_\Phi(p)$, by the computation in the proof of Lemma~\ref{lem:ol-symmetry}(c). So $\mathsf l=0$, $\Pi_\Phi(p)=O(|p|^2)$, and the marginal tensor of $\Pi_\Phi$ is $T^{\rm f}_{E^\infty}$. Its isotropic coefficient is $\beta^\infty_{E^\infty}$, by (b)(ii) in the form just proved.

(d) Let $X$ be a polymer of the infinite lattice. For $\tilde n$ large, depending on $X$, its projection to the torus of side $\tilde n$ is a polymer that is not wrapping, with the same cubes and the same tree length, since the projection is isometric on sets of diameter less than $\frac12\tilde n$, and $X$ is a lift of it. By \eqref{eq:torus-local}, $E^\infty(X,\cdot)$ is the uniform limit on $\mathcal U^\flat_1(X)$ of the functions $E[\tilde n]$ of the projections, read on the lift. So it is holomorphic, depends only on $B$ on the bonds of $X$, and is normalized, real and invariant under small real gauge transformations, as these functions are; and $|E^\infty(X,\cdot)|\le\mathsf K'e^{-\kappa_nd_j(X)}$. Property (G2) for $E^\infty$, the covariance $E^\infty(X+a,\tau_aU)=E^\infty(X,U)$ for $a\in M\Z^4$, is part of the definition of a torus-local family.

Let $\tilde n\in\mathfrak N_j$, let $Q$ be a grid cube of the torus and $\tilde Q$ a lift. Lifting with $Q\mapsto\tilde Q$ is a bijection from the non-wrapping polymers $X\supset Q$ of the torus onto the polymers $\tilde X\supset\tilde Q$ of the infinite lattice of diameter at most $\frac12\tilde n-8$. It preserves $|X|$. For such $X$ and $|F|\le1$, the function $t\mapsto E[\tilde n](X,V^X_{t,F})-E^\infty(\tilde X,V^{\tilde X}_{t,F})$, with the origin at the lift of a point of $X$, is holomorphic for $|t|<\delta_1^\flat/(dD_X)$ and bounded there by $Ce^{-c\tilde n/M}e^{-\kappa_nd_j(X)}$, by Lemma~\ref{lem:period-density}(a) and \eqref{eq:torus-local}. Cauchy's estimate gives
\[
\big|c^X_{E[\tilde n]}(F)-c^{\tilde X}_{E^\infty}(F)\big|\le9d^2M^2(\delta_1^\flat)^{-2}\big(1+d_j(X)\big)^2Ce^{-c\tilde n/M}e^{-\kappa_nd_j(X)} .
\]
The polymers $\tilde X\supset\tilde Q$ of larger diameter have $d_j(\tilde X)\ge\tilde n/(4M)$, as in Lemma~\ref{lem:torus}(b), and $|c^{\tilde X}_{E^\infty}(F)|\le9d^2M^2(\delta_1^\flat)^{-2}(1+d_j(\tilde X))^2\mathsf K'e^{-\kappa_nd_j(\tilde X)}$. Since $(1+t)^2e^{-(\kappa_n-\kappa_0)t}\le e^{-(\kappa_n-\kappa_0-2)t}$ and $\kappa_n-\kappa_0\ge3$, \eqref{eq:dM-facts} gives
\[
\Big|\sum_{X\supset Q}\frac{c^X_{E[\tilde n]}(F)}{|X|}-\sum_{\tilde X\supset\tilde Q}\frac{c^{\tilde X}_{E^\infty}(F)}{|\tilde X|}\Big|\le9d^2M^2(\delta_1^\flat)^{-2}K_0\big(Ce^{-c\tilde n/M}+\mathsf K'e^{-\tilde n/(4M)}\big),
\]
uniformly in $Q$. The second sum is $M^4\bar T_{E^\infty}(F)$, and $\bar T^{\rm vol}_{E[\tilde n]}(F)$ is $M^{-4}$ times the average over $Q$ of the first. So $|\bar T^{\rm vol}_{E[\tilde n]}(F)-\bar T_{E^\infty}(F)|\le9d^2K_0M^{-2}(\delta_1^\flat)^{-2}(C+\mathsf K')e^{-c'\tilde n/M}$ for $|F|\le1$, with $c'=\min(c,\frac14)$. Since $|F^{\mu\nu,a}|=1$ and there are $6\dim\mathfrak g$ such data, the bound on $\beta^{\rm vol}_{E[\tilde n]}-\beta^\infty_{E^\infty}$ follows, with $C'=18d^2K_0M^{-2}(\delta_1^\flat)^{-2}(C+\mathsf K')$. The kernels are compared in the same way: the torus kernel at two bonds is the sum, over the polymers that contain both, of the kernels of the second variations of $E[\tilde n](X,e^{\ii\eta A})$ ($2\kappa_X$ for $X$ not wrapping); the wrapping ones, by the same Cauchy estimate with $d_j(X)\ge\tilde n/(4M)$ (Lemma~\ref{lem:torus}(b)), and the polymers of the infinite lattice of diameter larger than $\frac12\tilde n-8$, contribute $O(e^{-c'\tilde n/M})$, and the others differ by $O(e^{-c\tilde n/M})$, by Cauchy's estimate, with constants depending on $\eta$. For the bound, by the same Cauchy estimate each of these kernels is at most $C_\eta\mathsf K'e^{-\kappa_nd_j(X)}$ at $b,b'$, and a polymer of the torus that contains $b$ and $b'$ has $d_j(X)\ge|m_b-m_{b'}|_{\tilde n}/M-2$, as in Lemma~\ref{lem:torus}(b). With \eqref{eq:dM-facts}, the torus kernel is at most $C_\eta K_0\mathsf K'e^{2(\kappa_n-\kappa_0)}e^{-(\kappa_n-\kappa_0)|m_b-m_{b'}|_{\tilde n}/M}$, uniformly in $\tilde n$.
\end{proof}

\paragraph{The step functions in infinite volume.} For $j\ge0$ let $\Phi_j(V')=\Gamma^{\rm step}_j(U_{j+1}(V'))$, a function of the configurations $V'$ of scale $j+1$ near $1$ (Section~\ref{sec:ol-composite}). Its marginal tensor is $T^{(1)}_{j+1}$.

\begin{lemma}[Infinite-volume limits of the step functions]\label{lem:ol-limits}
Let $G$ be compact, connected and simple, let $j\ge0$, and use the units of scale $j+1$, with $\eta=L^{-j-1}$. Let $E^\infty$ be the limit of the torus-local family of Lemma~\ref{lem:OL} with $k=j$, $\widetilde\Pi$ the kernel of its second variation (Section~\ref{subsec:ol-period-vacpol}), and $\bH_\infty$ the infinite-volume minimizer map of Proposition~\ref{prop:ol-IM} for the averaging to scale $j+1$. On the torus whose lattice of scale $j+1$ has side $\tilde n\in\mathfrak N_{j+1}$, the second variation of $\Phi_j$ at $V'=1$ is $\bH^{\sf T}\Pi_{\tilde n}\bH$, where $\Pi_{\tilde n}$ is the kernel of the second variation at $A=0$ of $A\mapsto\sum_{Y'}E'_{\rm ol}[\tilde n](Y',e^{\ii\eta A})$. As $\tilde n\to\infty$ in $\mathfrak N_{j+1}$, these second variations converge with uniform decay (Lemma~\ref{lem:ol-folding}(e)) to
\[
\Pi_{\Phi_j}=\bH_\infty^{\sf T}\widetilde\Pi\bH_\infty .
\]
In particular the hypothesis of Lemma~\ref{lem:ol-symmetry}(c), and the hypothesis of Corollary~\ref{cor:ol-cumulative} on the step functions, hold for $\Phi_j$.
\end{lemma}

\begin{proof}\hasdetail{lem:ol-limits}
Let $\tilde n\in\mathfrak N_{j+1}$. On this torus, for small-field $V'$ and $U'=U_{j+1}(V')$,
\[
\Gamma^{\rm step}_j(U')-\Gamma^{\rm step}_j(1)=\sum_X\mathfrak e(X,U')+\sum_c\log\mathfrak j_c(U')=\sum_{Y'}E'_{\rm ol}[\tilde n](Y',U'),
\]
by Proposition~\ref{prop:one-loop-polymer}(a) and Proposition~\ref{prop:jacobians}(b) (Proposition~\ref{prop:first-step}(c) at $j=0$), regrouped as in Definition~\ref{def:raw-new}. Let $V'=e^{\ii B}$ with $B$ small. As in the proof of Lemma~\ref{lem:ol-pullback}, $U_{j+1}(V')$ is a gauge transform of $e^{\ii\eta A_*(B)}$, with $A_*$ analytic and $A_*(B)=\bH B+O(|B|^2)$ (Corollary~\ref{cor:analytic}, Theorem~\ref{thm:variational}(c)); by the proof of Lemma~\ref{lem:ol-pullback}, $\bH$ is the minimizer map of Part~I. The function $A\mapsto\sum_{Y'}E'_{\rm ol}[\tilde n](Y',e^{\ii\eta A})$ is invariant under real gauge transformations, since its activities are, and its first variation at $A=0$ vanishes, since it is invariant under the constant gauge transformations and $G$ is semisimple. So only the linearization of $A_*$ enters, and the second variation of $\Phi_j$ at $V'=1$ is $\bH^{\sf T}\Pi_{\tilde n}\bH$.

By Lemma~\ref{lem:OL}, the family $(E'_{\rm ol}[\tilde n])_{\tilde n\in\mathfrak N_{j+1}}$ is torus-local, with limit $E^\infty$. By Lemma~\ref{lem:period-vacpol}(d), the kernels $\Pi_{\tilde n}$ converge to $\widetilde\Pi$ with uniform decay. By Proposition~\ref{prop:ol-IM}(a), the maps $\bH$ on the tori converge with uniform decay to $\bH_\infty$. By Lemma~\ref{lem:ol-folding}(e), applied to the transposes and to the two products, $\bH^{\sf T}\Pi_{\tilde n}\bH$ converges with uniform decay to $\bH_\infty^{\sf T}\widetilde\Pi\bH_\infty$.
\end{proof}

\subsection{The properties (F1)--(F3)}\label{subsec:ol-flow-F}

\begin{proposition}[The properties (F1)--(F3)]\label{prop:flow-F}
Let $G$ be compact, connected and simple, let the parameters and the torus be as in Section~\ref{sec:flow4}, and let $\gamma$ satisfy ($\Gamma$1)--($\Gamma$4) of Remark~\ref{rem:flow-constants}. Then (F1), (F2) and (F3) hold. In (F2),
\[
C_2=18d^2K_0M^{-2}(\delta_1^\flat)^{-2}\mathsf K,\qquad c_2=\tfrac14,
\]
the constants of Lemma~\ref{lem:period-vacpol}(d) for the one-loop parts. In (F3),
\[
C_3=C_\beta C_{\rm new}\epsilon_*\big(C_{\rm pure}+c_{\rm F}^{-2}\big),
\]
where $C_{\rm pure}$ is a constant with $r_t^{-2}\le C_{\rm pure}p_0^6$ for $p_0\ge1$, with $r_t$ of Proposition~\ref{prop:pure-sources} (such a constant is given in the proof of that proposition, where it is written $C_P$), and $c_{\rm F}$ is that of ($\Gamma$4); for the first step the constants of the first step are used, and $C_3$ is the larger of the two values.
\end{proposition}

\begin{proof}\hasdetail{prop:flow-F}
\emph{(F1).} By Theorem~\ref{thm:sf-flow}(a), at every $k$ with $g_1,\dots,g_k\le\gamma$, the densities of the flow have the generation form, with the recursion and $g_{j+1}\le2g_j$ of Theorem~\ref{thm:sf-flow}(e). The coefficients $\beta_{j+1}[\pi]$ are defined on $\Sigma_j$, with $|\beta_{j+1}|\le C_\beta\mathsf K$, by Proposition~\ref{prop:flow-continuity}(a), and they are continuous there by Proposition~\ref{prop:flow-continuity}(b); $p_0$ is continuous (Section~\ref{subsec:flow-parameters}).

\emph{(F3).} In outline: replace the coupling in the fluctuation by a complex parameter $t$. The coefficient $\phi(t)$ of the part $E'_\#$ of the raw generation beyond its one-loop part is then holomorphic and even in $t$, vanishes at $t=0$, and is bounded for $|t|<r_{\rm F}$, where $r_{\rm F}^{-2}\le(C_{\rm pure}+c_{\rm F}^{-2})p_0^6$; the maximum principle gives $|\phi(g)|\le Cg^2r_{\rm F}^{-2}$. In detail, fix $k$ and the sequence $g_0,\dots,g_k$, let $g=g_k$, $p_0=p_0(g)$, and let $\mathcal E_k$ be the old activities of step $k$ (Definition~\ref{def:old-activities}; none if $k=0$). As in Section~\ref{subsec:cl-order}, replace $g$ by a complex parameter $t$ in the fluctuation only, $Z(s)=tC_k^{1/2}(s)\xi$, with $t^{-2}$ in front of the interaction and gauge-fixing pieces, and keep the reference, the old activities and $d\mu^*$. Let $r_{\rm F}=\min(r_t,c_{\rm F}/p_0)$, with $r_t$ of Proposition~\ref{prop:pure-sources} and $c_{\rm F}$ of ($\Gamma$4); then $2g\le r_{\rm F}$, by ($\Gamma$2) and ($\Gamma$4).

Let $|t|<r_{\rm F}$ and $|\xi|_\infty\le p_0$. The pure pieces satisfy the bounds of Proposition~\ref{prop:localized-pieces} with $\epsilon_{\rm p}(t)\le\frac12\epsilon_*$ (Proposition~\ref{prop:pure-sources}). The proof of Lemma~\ref{lem:white-noise}, which needs here only $|t|C_Tp_0\le\frac12r_I\le r_L$ (the second entry of $r_t$), gives $|\hat Z(Z(s))|_\infty\le\zeta(t)=C_wL^dC_T|t|p_0$, and $\zeta(t)\le\frac12R_P$ and $\tau(t)=C_PL^d|t|p_0/\delta_1<\frac12$ (with the $C_P$ of Proposition~\ref{prop:old-activities}(d)), by the definition of $c_{\rm F}$ in ($\Gamma$4). So Proposition~\ref{prop:composed} applies to every old activity with $\zeta=\zeta(t)$, as in Proposition~\ref{prop:old-activities}(d), and the mixed pieces satisfy the bounds of Proposition~\ref{prop:localized-pieces} with $\epsilon_{\rm m}(t)=C_P\tau(t)C_{\rm old}\mathsf K$, where $\epsilon_{\rm m}(t)e^{c_\natural(\kappa_n-1)}\le\frac14\epsilon_*$. The pieces are analytic in $t$, jointly with the reference: the composed activities $f(X;s,Z)$ are analytic in $Z$, which is linear in $t$, and Lemma~\ref{lem:localization}(c) keeps the analyticity. So Corollary~\ref{cor:two-classes} applies for every such $t$, with its bounds at $\epsilon_{\rm p}(t)$, $\epsilon_{\rm m}(t)$, at the real critical reference $1$ and at the extensions of the points of the domains $\mathcal U(Y')$, through the restricted systems of Remark~\ref{rem:extension}. Let $E'_\#(t)$ be the family $E'_\#$ built from these outputs. By Theorem~\ref{thm:mayer}(c) and the proof of Lemma~\ref{lem:raw-new}(a), (b), $E'_\#(t)(Y',U)$ is holomorphic in $(t,U|_{Y'})$, normalized and gauge invariant, and
\[
\|E'_\#(t)\|^{\rm raw}_{k+1}\le C_{\rm new}\big(\epsilon_{\rm p}(t)+\epsilon_{\rm m}(t)\big)\le C_{\rm new}\epsilon_*\qquad(|t|<r_{\rm F}).
\]

\emph{Parity.} Every piece depends on $(t,\xi)$ only through $t^{-2}$ and $Z(s)$: the pure pieces by the proof of Proposition~\ref{prop:pure-sources}, and the mixed pieces since $f(X;s,Z(s))$ depends only on $Z(s)$. So they are even under $(t,\xi)\mapsto(-t,-\xi)$, and $E'_\#(t)$ is even in $t$, by Theorem~\ref{thm:mayer}(d). At $t=0$ every piece vanishes: the pure pieces by Proposition~\ref{prop:pure-sources}, and the composed activities since $Z=0$ gives $\hat Z(0)=0$, $A_*(s;0)=0$ (Theorem~\ref{thm:weak-bg}(a)), and so $f(X;s,0)=\mathcal E_k(X,U')-\mathcal E_k(X,U')=0$. Then every $K(Z)=e^{E(Z)}-1$ in the proof of Theorem~\ref{thm:mayer} vanishes, and $E'_\#(0)=0$.

\emph{Cauchy.} Put $\phi(t)=\beta^{\rm vol}_{E'_\#(t)}$. Each $c^X_{E'_\#(t)}(F)$ is a Cauchy integral, in the variable of $V^X_{\cdot,F}$, of a function holomorphic in $t$, and the series \eqref{eq:beta-vol} converges uniformly in $t$, by the bound of Lemma~\ref{lem:period-density}(c). So $\phi$ is holomorphic on $|t|<r_{\rm F}$, even, with $\phi(0)=0$, and $|\phi(t)|\le C_\beta C_{\rm new}\epsilon_*$ there, by \eqref{eq:beta-vol-bound}. Hence $\phi(t)/t^2$ is holomorphic, and the maximum principle gives $|\phi(g)|\le C_\beta C_{\rm new}\epsilon_*(g/r_{\rm F})^2$. Moreover $r_{\rm F}^{-2}\le r_t^{-2}+p_0^2c_{\rm F}^{-2}\le(C_{\rm pure}+c_{\rm F}^{-2})p_0^6$, since $r_t^{-2}\le C_{\rm pure}p_0^6$ and $p_0\ge1$. At $t=g$ the outputs are those of the step, so $E'_\#(g)=E'-E'_{\rm ol}$ and $\beta_{k+1}-\beta^{\rm ol}_{k+1}=\phi(g)$. This is (F3). At $k=0$ there is no mixed class, $r_{\rm F}=r_t$, and the constants are those of the first step.

\emph{(F2).} Fix $k<K$. By Lemma~\ref{lem:OL}, the one-loop parts $E'_{\rm ol}[\tilde n]$, $\tilde n\in\mathfrak N_{k+1}$, form a torus-local family, real at real configurations, with $\|E'_{\rm ol}[\tilde n]\|\le\mathsf K$, for which the left side of \eqref{eq:torus-local} vanishes; let $E^\infty$ be its limit. The torus of the flow has its side $n_{k+1}$ at scale $k+1$ in $\mathfrak N_{k+1}$. By Lemma~\ref{lem:period-vacpol}(d),
\[
\big|\beta^{\rm ol}_{k+1}-\beta^\infty_{E^\infty}\big|\le C_2e^{-c_2n_{k+1}/M}.
\]
Indeed, by Lemma~\ref{lem:OL} the left side of \eqref{eq:torus-local} vanishes. So, in the proof of Lemma~\ref{lem:period-vacpol}(d), the polymers of diameter at most $\frac12\tilde n-8$ contribute nothing to the difference, and only the polymers of the infinite lattice of larger diameter remain. With $\mathsf K'\le\mathsf K$ they give the bound with $C_2=18d^2K_0M^{-2}(\delta_1^\flat)^{-2}\mathsf K$ and $c_2=\frac14$.

By Lemma~\ref{lem:ol-limits} with $j=k$, the second variations at $V'=1$ of $V'\mapsto\Gamma^{\rm step}_k(U_{k+1}(V'))$ on the tori of the sides in $\mathfrak N_{k+1}$ converge with uniform decay to $\Pi_\Phi=\bH_\infty^{\sf T}\widetilde\Pi\bH_\infty$. This is the infinite-volume second variation of Definition~\ref{def:marginal-tensor}, so $T^{(1)}_{k+1}$ is the marginal tensor of $\Pi_\Phi$. By Lemma~\ref{lem:ol-symmetry}(b), these functions are invariant under the translations and under $B_4$ about the coarse points, so their second variations at $1$ are, and so is $\Pi_\Phi$. Lemma~\ref{lem:period-vacpol}(c) gives $T^{(1)}_{k+1}=T^{\rm f}_{E^\infty}$, with isotropic coefficient $\beta^\infty_{E^\infty}$. By Lemma~\ref{lem:ol-symmetry}(c), $T^{(1)}_{k+1}=\beta^{(1)}_{k+1}(\delta_{\mu\nu}|p|^2-p_\mu p_\nu)$, so $\beta^{(1)}_{k+1}=\beta^\infty_{E^\infty}$. With the bound above this is (F2). The constants $C_2$, $c_2$ depend only on $d$, $M$, $\kappa_0$ (through $K_0$), $\delta_1^\flat$ and $\mathsf K$, not on $k$, $K$ or the torus.
\end{proof}

\begin{remark}[On (F1)--(F3)]\label{rem:flow-F}
\begin{enumerate}[label=(\alph*)]
\item \emph{The old generations.} The old generations enter the step only through the old activities $\mathcal E_k$, in which every counterterm $-\beta_jS_\Box$ is grouped with the small polymers of its generation in the cube $\Box$ (Definition~\ref{def:old-activities}). So their marginal parts are cancelled, for the generations with $k-j\ge m_0$ (Proposition~\ref{prop:extraction}), or grouped with their counterterm in bounded old activities, for the recent ones (Proposition~\ref{prop:old-activities}), before the new generation is formed. The old activities do not enter $E'_{\rm ol}$, and they enter $\beta_{k+1}$ only through the mixed class in $E'_\#$, which (F3) bounds by $O(g_k^2p_{0,k}^6)$.
\item \emph{Sizes.} Without the parity, Corollary~\ref{cor:two-classes} would give only $O(\epsilon_{\rm p}+\epsilon_{\rm m})=O(gp_0^3)$ in (F3), which is not summable along the flow. The errors of (F2) are summable over the scales: $\sum_{j<K}e^{-c_2n_{j+1}/M}\le\sum_{m\ge0}e^{-c_2L^m}$, since $n_{j+1}=L^{K-j-1}n_K$ and $n_K\ge M$.
\end{enumerate}
\end{remark}

The next proposition uses Section~\ref{sec:gen-locality}. That section uses Parts~IV--X, Lemma~\ref{lem:ol-folding}, and Sections~\ref{subsec:ol-flow-coeff} and~\ref{subsec:ol-period-vacpol} with Lemmas~\ref{lem:OL} and~\ref{lem:period-vacpol}, but none of the results below.

\begin{proposition}[The coefficients and Ba\l aban's $\beta$]\label{prop:beta-identification}
Let $G$ be compact, connected and simple, let $j\ge1$, let $\gamma$ satisfy ($\Gamma$1)--($\Gamma$3), and let $(g_0,\dots,g_{j-1})\in\Sigma_{j-1}$. For $\tilde n\in\mathfrak N_j$ let $E'[\tilde n]$, $E^{(j)}[\tilde n]=E'[\tilde n]^{\rm sym}$ and $\beta_j[\tilde n]=\beta_{E^{(j)}[\tilde n]}$ be the raw generation, the generation and the coefficient of scale $j$ constructed with these couplings on the torus whose lattice of scale $j$ has side $\tilde n$ (Proposition~\ref{prop:flow-continuity}, with $K=j$). Then:
\begin{enumerate}[label=(\alph*)]
\item the second variations at $V=1$ of $\Phi_{E^{(j)}[\tilde n]}$ converge as $\tilde n\to\infty$ in $\mathfrak N_j$, with uniform decay (Lemma~\ref{lem:ol-folding}(e)), to a kernel $\Pi_\Phi$ with exponential decay, invariant under the translations and under $B_4$, as Definition~\ref{def:marginal-tensor} assumes;
\item its marginal tensor is $\beta^{\rm B}_j(\delta_{\mu\nu}|p|^2-p_\mu p_\nu)$, where $\beta^{\rm B}_j=-\partial_{p_\mu}\partial_{p_\nu}\Pi_\Phi(p)_{\mu\nu}|_{p=0}$, $\mu\ne\nu$, is the coefficient that Ba\l aban's definition \cite[(1.20)--(1.22)]{BalYM4I} assigns to the family $(E^{(j)}[\tilde n])_{\tilde n\in\mathfrak N_j}$ built with these couplings, normalized relative to $S$ (Definition~\ref{def:marginal-tensor}). It equals $\beta^\infty_j$ of Theorem~\ref{thm:gen-locality}, and it depends only on $(g_0,\dots,g_{j-1})$. Ba\l aban's own generation is a different construction, which depends only on $g_{j-1}$;
\item $\beta_j[\tilde n]\to\beta^{\rm B}_j$ as $\tilde n\to\infty$ in $\mathfrak N_j$.
\end{enumerate}
\end{proposition}

\begin{proof}\hasdetail{prop:beta-identification}
Let $h'$ be a threshold at which $E'[\tilde n]$ is a raw generation (Theorem~\ref{thm:gen-step}(b) with $k=j-1$; for $j=1$, Proposition~\ref{prop:gen-one}). By Lemma~\ref{lem:symmetrization}(b), $\Phi_{E^{(j)}[\tilde n]}=\Phi_{E'[\tilde n]}$ on $\mathfrak U_j(h')$, which contains a neighbourhood of $V=1$, so their second variations at $V=1$ agree. As in the proof of Lemma~\ref{lem:ol-limits}, with $\Phi_{E'[\tilde n]}(V)=\sum_YE'[\tilde n](Y,U_j(V))$ in place of $\Phi_j$, the second variation is $\bH^{\sf T}\Pi_{\tilde n}\bH$, where $\Pi_{\tilde n}$ is the kernel of the second variation at $A=0$ of $A\mapsto\sum_YE'[\tilde n](Y,e^{\ii\eta A})$. By Theorem~\ref{thm:gen-locality}(a), the collection $(E'[\tilde n])$ converges on the flat domains $\mathcal U_1$ and has the covariance of Lemma~\ref{lem:gl-limits}(b); its members are as in Lemma~\ref{lem:gl-limits}(c), with $\|E'[\tilde n]\|\le\|E'[\tilde n]\|^{\rm raw}_j\le\mathsf K$. So, by Lemma~\ref{lem:gl-limits}(c), the restriction of its limit $E'^{(j),\infty}$ to the domains $\mathcal U^\flat_1$ is a family of flat activities with (G2), and the kernels $\Pi_{\tilde n}$ converge with uniform decay to its kernel $\widetilde\Pi$. By Proposition~\ref{prop:ol-IM}(a), the maps $\bH$ on the tori converge with uniform decay to $\bH_\infty$. By Lemma~\ref{lem:ol-folding}(e), applied to the transposes and to the two products, the second variations converge with uniform decay to $\Pi_\Phi=\bH_\infty^{\sf T}\widetilde\Pi\bH_\infty$. By (G3) for $E'[\tilde n]$ (Theorem~\ref{thm:gen-step}(b), and Proposition~\ref{prop:gen-one} for $j=1$), $\Phi_{E'[\tilde n]}$ is invariant under $\mathfrak S_j$ near $V=1$, so the kernels and $\Pi_\Phi$ are invariant under the translations and under $B_4$. By Lemma~\ref{lem:period-vacpol}(c), the marginal tensor is $T^{\rm f}_{E'^{(j),\infty}}$, and by the tensor algebra in the proof of Lemma~\ref{lem:ol-symmetry}(c), with the transversality of $\Pi_\Phi$ from Remark~\ref{rem:marginal-orders} ($\Phi_{E'[\tilde n]}$ is invariant under real gauge transformations, by Lemma~\ref{lem:raw-new}(a) and the paragraph before Definition~\ref{def:generation}, and its first variation at $1$ vanishes, since $G$ is semisimple), it is isotropic, hence equal to $\beta^\infty_{E'^{(j),\infty}}(\delta_{\mu\nu}|p|^2-p_\mu p_\nu)$. For $\mu\ne\nu$ the entry $\mu\nu$ of this tensor is $-\beta^\infty_{E'^{(j),\infty}}p_\mu p_\nu$, and $\Pi_\Phi(p)$ differs from the tensor by $O(|p|^3)$; so $\beta^{\rm B}_j=\beta^\infty_{E'^{(j),\infty}}=\beta^\infty_j$. The limit depends only on the couplings, by Theorem~\ref{thm:gen-locality}. Finally (c) is Theorem~\ref{thm:gen-locality}(c).
\end{proof}

The comparison is made with the counterterm normalized by $S$ itself. In the trace pairing of Section~\ref{sec:action}, the marginal tensor of $S$ has isotropic coefficient $1$ (proof of Lemma~\ref{lem:ol-classical}), so $\beta^{\rm B}_j$ is also the ratio of the coefficient of $\Phi_{E^{(j)}}$ to that of $S$; in this form the comparison does not depend on the trace convention. Theorem~\ref{thm:gen-locality} plays, for the generations of this note, the role of the existence of the limit that Ba\l aban asserts for (1.20) ``by the localized representation (1.7)''. It is qualitative: no rate is proved, and the convergence is not shown to be uniform in the couplings (Remark~\ref{rem:gl-rate}).

\subsection{The ultraviolet limit}\label{subsec:ol-uv-limit}

\begin{theorem}[The ultraviolet limit]\label{thm:ol-uv}
Let $G$ be compact, connected and simple, let the parameters and the torus be as in Section~\ref{sec:flow4}, with $n_K\ge n_{\min}$, and let $p_0(g)=(\log g^{-1})^q$. Let $\gamma$ satisfy ($\Gamma$1)--($\Gamma$4), and moreover $\gamma\le e^{-3q}$ and $C_3\gamma^2p_0(\gamma)^6\le\frac12\beta_\infty$, with $\beta_\infty$ of Theorem~\ref{thm:ol-cumulative} and $C_3$ of Proposition~\ref{prop:flow-F}. There are $g_{\rm uv}>0$ and $C$, independent of $K$ and of the torus, such that for $0<g\le g_{\rm uv}$ and every $K$ there is a bare coupling $g_0$, depending on $K$, $g$ and the torus, for which the flow has $g_k\in(0,\gamma]$ for $k\le K$ and $g_K=g$, and, for every $k\le K$,
\[
\Big|\frac1{g_k^2}-\frac1{g^2}-(K-k)\beta_\infty\Big|\le C\Big(1+\sum_{j=k}^{K-1}g_j^2p_0(g_j)^6\Big)\le C\big(\log(K-k+2)\big)^{6q+1}.
\]
In particular, there is $m_1$, independent of $K$, $g$ and the torus, such that $\frac12(K-k)\beta_\infty\le g_k^{-2}-g^{-2}\le2(K-k)\beta_\infty$ for $K-k\ge m_1$, and $g_k^{-2}-g^{-2}\le C(K-k)$ for all $k\le K$.
\end{theorem}

Ba\l aban's Theorem~2 \cite{BalYM4I}, stated for $d=4$ and $G=SU(2)$, asserts that for $g$ small there is a bare coupling for which his couplings satisfy $g_k\in\,]0,\gamma]$, $g_K=g$ and, with constants $0<\beta\le\beta'$,
\[
\frac1{g^2}+\beta\log(L^k\varepsilon)^{-1}\le\frac1{g_k^2}\le\frac1{g^2}+\beta'\log(L^k\varepsilon)^{-1},
\]
his (0.31) \cite[p.~259]{BalYM4I}. With $\varepsilon=L^{-K}$ the finest spacing in the units of the last scale, $\log(L^k\varepsilon)^{-1}=(K-k)\log L$. So the last statement of Theorem~\ref{thm:ol-uv} is the analogue of (0.31) for the couplings of the small-field flow, with $\beta=\beta_\infty/(2\log L)$ and $\beta'=2\beta_\infty/\log L$, which satisfy $0<\beta\le\beta'$, for $K-k\ge m_1$, and its upper half for all $k$, with $\beta'=\max(2\beta_\infty,C_\beta\mathsf K)/\log L$. The lower half for the last steps, $1\le K-k<m_1$, would need positivity of the single coefficients $\beta_{j+1}$, $K-m_1\le j<K$. Theorem~\ref{thm:ol-cumulative} controls only the cumulative sums $c_k=\sum_{j<k}\beta^{(1)}_{j+1}$, and it does not show that a single coefficient $\beta^{(1)}_{j+1}=c_{j+1}-c_j=\beta_\infty+O(1)$ is positive; nor is the error of (F2) small at $j=K-1$ unless $n_K$ is large. There the lower bound holds with an additive constant, $g_k^{-2}-g^{-2}\ge\frac12(K-k)\beta_\infty-C$ (by \eqref{eq:uv-lower} below). A possible route to the exact form is described in Remark~\ref{rem:ol-remains}(c).

\begin{proof}\hasdetail{thm:ol-uv}
\emph{Cumulative coefficients.} Put $c_k=\sum_{j<k}\beta^{(1)}_{j+1}$. By Theorem~\ref{thm:ol-cumulative}, $c_k=k\beta_\infty+r_k$ with $|r_k|\le C$. Here $C$ does not depend on $K$ or on the torus, since the tensors $T^{(1)}_{j+1}$ are infinite-volume quantities (Definition~\ref{def:marginal-tensor}). In particular each $\beta^{(1)}_{j+1}$ is bounded, though it may be negative for small $j$; this does no harm.

\emph{Localized error.} For $0\le k\le K$ and couplings $g_k,\dots,g_{K-1}$ in $(0,\gamma]$, (F2) and (F3) give
\[
\sum_{j=k}^{K-1}\big|\beta_{j+1}-\beta^{(1)}_{j+1}\big|\le C_3\sum_{j=k}^{K-1}g_j^2p_0(g_j)^6+C_4\le\tfrac12(K-k)\beta_\infty+C_4,
\]
with $C_4=C_2\sum_{m\ge0}e^{-c_2L^m}$ (Remark~\ref{rem:flow-F}(b)). The second inequality uses that $g\mapsto g^2p_0(g)^6$ increases on $(0,e^{-3q}]$, and $C_3\gamma^2p_0(\gamma)^6\le\frac12\beta_\infty$. Hence, along any flow that stays in $(0,\gamma]$ from scale $k$ to scale $K$,
\begin{equation}\label{eq:uv-lower}
\frac1{g_k^2}=\frac1{g_K^2}+\sum_{j=k}^{K-1}\beta_{j+1}\ge\frac1{g_K^2}+(c_K-c_k)-\tfrac12(K-k)\beta_\infty-C_4\ge\frac1{g_K^2}+\tfrac12(K-k)\beta_\infty-C_5,
\end{equation}
with $C_5=2C+C_4$.

\emph{Exit argument.} By Corollary~\ref{cor:bare-continuity}, the set $\mathsf B_K$ of the $g_0\in(0,\gamma]$ for which $g_1,\dots,g_K\in(0,\gamma]$ is closed in $(0,\gamma]$ and contains an interval $(0,\bar g]$, the map $g_0\mapsto g_K$ is continuous on it, and $g_K\to0$ as $g_0\to0$, since $|g_K^{-2}-g_0^{-2}|\le KC_\beta\mathsf K$. Let $I=(0,g_0^*]$ be the connected component of $\mathsf B_K$ that contains $(0,\bar g]$; it is closed in $(0,\gamma]$, since $\mathsf B_K$ is. Some $g_j$, $j\le K$, equals $\gamma$ at $g_0^*$. If $g_0^*=\gamma$, this holds with $j=0$. Otherwise there are $g_0\downarrow g_0^*$ outside $\mathsf B_K$, since $I$ is a component of the closed set $\mathsf B_K$. At each of them a first coupling $g_j$, $1\le j\le K$, exceeds $\gamma$, and $g_0\in\mathsf B_{j-1}$; the index $j$ is the same along a subsequence. Since $g_j$ is continuous on $\mathsf B_{j-1}\ni g_0^*$ (Corollary~\ref{cor:bare-continuity}(b)), $g_j(g_0^*)\ge\gamma$, hence $g_j(g_0^*)=\gamma$, since $g_0^*\in\mathsf B_K$. The localized estimate \eqref{eq:uv-lower} from $j$ to $K$ then gives
\[
\frac1{g_K^2}\le\frac1{\gamma^2}+C_5 ,
\]
so $g_K\ge\gamma(1+C_5\gamma^2)^{-1/2}>g$ for $g\le g_{\rm uv}$ small. By the intermediate value theorem, some $g_0\in I$ gives $g_K=g$.

\emph{The bound.} For this flow and every $k\le K$,
\[
\frac1{g_k^2}-\frac1{g^2}=\sum_{j=k}^{K-1}\beta_{j+1}=(K-k)\beta_\infty+r_K-r_k+\sum_{j=k}^{K-1}\big(\beta_{j+1}-\beta^{(1)}_{j+1}\big),
\]
and the last sum is at most $C_3\sum_{j=k}^{K-1}g_j^2p_0(g_j)^6+C_4$ in absolute value, by (F2) and (F3). This gives the first inequality. For the second, \eqref{eq:uv-lower} gives $g_j^2\le4/((K-j)\beta_\infty)$ for $K-j\ge4C_5/\beta_\infty$, and then, by the monotonicity, $g_j^2p_0(g_j)^6\le C(\log(K-j+2))^{6q}/(K-j)$. The other terms are at most $\gamma^2p_0(\gamma)^6$, and there are at most $4C_5/\beta_\infty$ of them. Summing over $j$ gives $C(\log(K-k+2))^{6q+1}$. The last statement follows, since $(\log(m+2))^{6q+1}/m\to0$ as $m\to\infty$, and since $|\beta_{j+1}|\le C_\beta\mathsf K$ for the upper bound when $K-k<m_1$.
\end{proof}

\begin{remark}[What remains]\label{rem:ol-remains}
\begin{enumerate}[label=(\alph*)]
\item \emph{The flow.} The flow is Part~X, and the properties (F1)--(F3) it gives the ultraviolet limit are proved in Proposition~\ref{prop:flow-F}. Property (F2), the comparison of the one-loop part of $\beta_{j+1}$, read off on the torus, with the coefficient of Part~XI, read off in infinite volume, rests on the torus independence of the one-loop pieces and of the Jacobian terms (Proposition~\ref{prop:one-loop-local}, Lemma~\ref{lem:OL}), on the infinite-volume minimizer (Proposition~\ref{prop:ol-IM}) and on the infinite-volume limits of the second variations (Lemma~\ref{lem:ol-limits}). The comparison with Ba\l aban's (0.31) is the paragraph after Theorem~\ref{thm:ol-uv}. A smooth small-field cutoff, as in \cite[(2.9)]{BalYM4I}, should give $O(g^2)$ in (F3) and the bound $O(\log(K-k+2))$. Theorem~\ref{thm:ol-uv} concerns the small-field model (Remark~\ref{rem:sf-flow}(d)).
\item \emph{The linear term.} With Ba\l aban's symmetric averaging, adopted in Part~IV, the functions $\Gamma^{(1)}_k$ and $\Gamma^{\rm step}_j\circ U_{j+1}$ are invariant under the translations and under $B_4$ (Lemma~\ref{lem:ol-symmetry}). So their vacuum polarizations have no term linear in $p$, and their marginal tensors are isotropic. The symmetry needs only the averaging: Lemma~\ref{lem:ol-symmetry}(a), and the route to~(b) through~(a), use no covariance property of the gauge fixing of Part~VII; the direct proof of~(b) uses only the covariance of its quadratic part, which is that of the averaged contour variables.
\item \emph{What remains.}
  \begin{itemize}
  \item \emph{A rate for generation locality.} Generation locality is proved in Section~\ref{sec:gen-locality}, qualitatively, so Proposition~\ref{prop:beta-identification} holds without hypothesis. Not proved: a rate in $\beta_j[\tilde n]\to\beta^{\rm B}_j$, and uniformity in the couplings (Remark~\ref{rem:gl-rate}). Proposition~\ref{prop:beta-identification} is not used by Theorem~\ref{thm:ol-uv}.
  \item \emph{The last steps.} For the last $m_1$ steps the lower half of the analogue of Ba\l aban's (0.31) holds only up to an additive constant (the paragraph after Theorem~\ref{thm:ol-uv}). The exact form would need the single coefficients $\beta_{j+1}$, $K-m_1\le j<K$, to be positive. A possible route, for $K$ large, is the convergence to $\beta_\infty$ of $\beta^{(1)}_{j+1}=c_{j+1}-c_j$, where $c_k$ is the coefficient of $T_{\Gamma^{(1)}_k}$ (Theorem~\ref{thm:ol-cumulative}). By Lemma~\ref{lem:ol-coarse-map} and Proposition~\ref{prop:ol-reduction}, $c_k$ is $\beta_{\rm F}$ at $\eta=L^{-k}$ plus the coefficient $r^{\rm fine}_k$ of the remainder $T_{\rm fine}-T_{\rm F}$, which is bounded uniformly in $k$. By Theorem~\ref{thm:ol-log}, $\beta_{\rm F}-k\beta_\infty$ converges as $k\to\infty$, so the differences of the first terms converge to $\beta_\infty$. What is missing is the convergence of $r^{\rm fine}_k$ as $k\to\infty$, a continuum limit at $U=1$ of explicit operators. In addition, the errors of (F2) and (F3) at the last steps must be small compared with $\beta_\infty$, which needs a lower bound on $n_K$ and a condition on $\gamma$. This is not carried out here.
  \item \emph{The large-field stage}, which is not treated (Remark~\ref{rem:sf-flow}(d)).
  \end{itemize}
\item \emph{Numerical checks.} Theorem~\ref{thm:ol-chart} (Remark~\ref{rem:ol-numerical}), the constant of Remark~\ref{rem:one-loop-intrinsic} and the invariance of Lemma~\ref{lem:ol-symmetry}, all three in $d=2$; the coefficient $\mathsf b$, its sign and the lattice constant of Dashen and Gross, for $SU(2)$ (Remark~\ref{rem:ol-f-numerical}); and the continuum tensor and, for $SU(2)$, the scaling of the integrand (Remark~\ref{rem:ol-integrand-numerical}) were checked numerically. No proof uses these checks. Appendix~\ref{app:numerics} describes them.
  The scaling check of the lattice integrand (Remark~\ref{rem:ol-integrand-numerical}(ii)) and the earlier checks are independent of the proofs; the spherical average (Remark~\ref{rem:ol-integrand-numerical}(i)) recomputes that of Lemma~\ref{lem:ol-classical} in exact arithmetic from the vertices of Lemma~\ref{lem:ol-naive-limit}. They test the formula of Theorem~\ref{thm:ol-chart}, the invariance of Lemma~\ref{lem:ol-symmetry}, and the lattice vertices, the joint expansion and the normalization of Section~\ref{sec:ol-feynman}. They say nothing about the remainders of Proposition~\ref{prop:ol-reduction}.
\end{enumerate}
\end{remark}
